\section{Original maximal order and tangent directions}
\label{sec:maximal-contact}

\subsection{Marked ideals of maximal order and tangent directions}

In this subsection \((\mathfrak a,\mu)\) is a marked ideal on a complex manifold \(X\) with \(\mathfrak a\) coherent and \(\mu\ge1\),
\(\mathcal D^i\) are the derivative ideals of Definition~\ref{hir:derivative-ideal}, and
\(\operatorname{cosupp}(\mathfrak a,\mu)=\{x:\operatorname{ord}_x\mathfrak a\ge\mu\}\) (Definition~\ref{hir:cosupport}).

\begin{definition}[Maximal order]
\label{mc:maximal-order}
\uses{hir:derivative-ideal,hir:cosupport}
\source{wlodarczyk05-hironaka}
\lean{Hironaka.MaximalContact.IsMaximalOrder}
\leanok
\((\mathfrak a,\mu)\) is \emph{of maximal order} if \(\mathcal D^\mu(\mathfrak a)=\mathcal O_X\), i.e.\ \(1\in\mathcal D^\mu(\mathfrak a)_x\) for every
\(x\in X\).
\end{definition}
By Lemma~\ref{hir:order-taylor} this says \(\operatorname{ord}_x\mathfrak a\le\mu\) for every \(x\), Włodarczyk's definition.

\begin{lemma}[Derivatives of a marked ideal of maximal order]
\label{mc:maximal-order-derivative}
\uses{mc:maximal-order}
\source{wlodarczyk05-hironaka}
\lean{Hironaka.MaximalContact.iterDerivIdeal_iterDerivIdeal,Hironaka.MaximalContact.IsMaximalOrder.iterDerivIdeal}
\leanok
If \((\mathfrak a,\mu)\) is of maximal order and \(0\le i\le\mu\), then \((\mathcal D^i(\mathfrak a),\mu-i)\) is of maximal order.
\end{lemma}
\begin{proof}
\leanok
\(\mathcal D^{\mu-i}(\mathcal D^i(\mathfrak a))=\mathcal D^\mu(\mathfrak a)=\mathcal O_X\), by the definition \(\mathcal D^{j+1}=\mathcal D\circ\mathcal D^j\) and induction on
\(\mu-i\).
\end{proof}

\begin{definition}[Tangent direction]
\label{mc:tangent-direction}
\uses{hir:derivative-ideal}
\source{wlodarczyk05-hironaka}
\lean{Hironaka.MaximalContact.IsTangentDirection}
\leanok
Let \(U\subset X\) be open. A \emph{tangent direction} of \((\mathfrak a,\mu)\) on \(U\) is a function \(u\) whose germ at every
point of \(U\) lies in \(\mathcal D^{\mu-1}(\mathfrak a)\) and which has multiplicity one: \(du(x)\ne0\) (equivalently, some partial
derivative \(\partial^e_vu(x)\) in some chart \(e\) at \(x\) is nonzero) for every \(x\in U\) with \(u(x)=0\).
\end{definition}
By the implicit function theorem the zero set \(V(u)\) of a tangent direction is a closed submanifold of \(U\) of codimension
one (or empty), a \emph{hypersurface of maximal contact}.

\begin{lemma}[The cosupport lies on every tangent direction]
\label{mc:tangent-direction-support}
\uses{mc:tangent-direction,hir:giraud-support}
\source{wlodarczyk05-hironaka}
\lean{Hironaka.MaximalContact.cosupport_add_subset_cosupport_iterDerivIdeal,Hironaka.MaximalContact.apply_eq_zero_of_mem_cosupport,Hironaka.MaximalContact.IsTangentDirection.cosupport_inter_subset}
\leanok
If \(u\) is a tangent direction of \((\mathfrak a,\mu)\) on \(U\), then \(\operatorname{cosupp}(\mathfrak a,\mu)\cap U\subset V(u)\).
\end{lemma}
\begin{proof}
\leanok
By Lemma~\ref{hir:giraud-support} (the inclusion \(\operatorname{cosupp}(\mathfrak a,j+1)\subset\operatorname{cosupp}(\mathcal D(\mathfrak a),j)\), applied
\(\mu-1\) times), \(x\in\operatorname{cosupp}(\mathfrak a,\mu)\) implies \(x\in\operatorname{cosupp}(\mathcal D^{\mu-1}(\mathfrak a),1)\), i.e.\ every germ of
\(\mathcal D^{\mu-1}(\mathfrak a)_x\) vanishes at \(x\); in particular \(u(x)=0\).
\end{proof}

\begin{lemma}[Existence of tangent directions]
\label{mc:tangent-direction-exists}
\uses{mc:maximal-order,mc:tangent-direction,mc:tangent-direction-support,hir:derivative-ideal-generators}
\source{wlodarczyk05-hironaka}
\lean{Hironaka.MaximalContact.exists_isTangentDirection}
\leanok
If \((\mathfrak a,\mu)\) is of maximal order and \(x\in\operatorname{cosupp}(\mathfrak a,\mu)\), then \((\mathfrak a,\mu)\) has a tangent direction on
some open neighbourhood of \(x\).
\end{lemma}
\begin{proof}
\leanok
\(\mathcal D^{\mu-1}(\mathfrak a)\) is coherent (Lemma~\ref{hir:derivative-ideal-generators}); let \(f_1,\dots,f_m\) generate it on an open
neighbourhood \(U_0\) of \(x\) inside the domain of a chart \(e\). By the same lemma, \(\mathcal D(\mathcal D^{\mu-1}(\mathfrak a))=\mathcal D^\mu(\mathfrak a)\) is
generated near \(x\) by the \(f_i\) and the \(\partial^e_{e_j}f_i\); since it contains \(1\), not all of these vanish at \(x\). As in the proof
of Lemma~\ref{mc:tangent-direction-support}, every germ of \(\mathcal D^{\mu-1}(\mathfrak a)_x\) vanishes at \(x\), so \(f_i(x)=0\) for all \(i\), and
\(\partial^e_{e_j}f_i(x)\ne0\) for some \(i,j\). Then \(df_i(x)\ne0\), hence \(df_i\ne0\) on an open neighbourhood \(U\subset U_0\) of \(x\), and
\(u=f_i|_U\) is a tangent direction on \(U\).
\end{proof}

\subsection{Derivative ideals under a blowup}

In this subsection \(\pi\colon Y\to X\) is the blowup of the complex manifold \(X\) along \(C\)
(Definition~\ref{ana:blowups}), \(D\) its exceptional divisor, so that \(\mathcal O_Y(-qD)=\pi^*(\mathcal I_C^q)\)
(Lemma~\ref{bcalc:exceptional-divisor}), and \(\sigma^c(\mathfrak a,\mu)=(\pi^*\mathfrak a:\mathcal O_Y(-\mu D))\) is the
controlled transform (Definition~\ref{bcalc:controlled-transform-def}). Marked ideals are taken with \(N\) the whole
manifold. For a chart \(e\) and \(v\in\mathbf C^n\), \(\partial^e_v\) is the partial derivative of
Definition~\ref{hir:derivative-ideal}; \(e_z\) is a chart at \(z\) and \(e_1,\dots,e_n\) the standard basis.

\begin{lemma}[Villamayor: maximal order survives admissible blowups]
\label{mc:maximal-order-transform}
\uses{mc:inclusions,mc:maximal-order,hir:hadamard-bridge}
\source{wlodarczyk05-hironaka}
\lean{Hironaka.MaximalContact.IsMaximalOrder.controlledTransform}\upstreamlean{TranscendentalNakamaye.Pending.MaximalContact.IsMaximalOrder.controlledTransform_of_subset_cosupport}
\mixedleanok
Let \((\mathfrak a,\mu)\) be a coherent marked ideal of maximal order with \(\mathfrak a\subset\mathcal I_C^\mu\) (for instance
\(C\subset\operatorname{cosupp}(\mathfrak a,\mu)\)). Then \((\sigma^c(\mathfrak a,\mu),\mu)\) is of maximal order.
\end{lemma}
\begin{proof}
\mixedleanok
By Lemma~\ref{mc:inclusions} with \(r=\mu\) and mark \(0\), \(\sigma^c(\mathcal D^\mu(\mathfrak a),0)\subset\mathcal D^\mu(\sigma^c(\mathfrak a,\mu))\). Since
\(1\in\mathcal D^\mu(\mathfrak a)\) everywhere, \(1=1\circ\pi\in\pi^*\mathcal D^\mu(\mathfrak a)=\sigma^c(\mathcal D^\mu(\mathfrak a),0)\).
\end{proof}
\begin{lemma}[The first centre meeting the smooth part]
\label{mc:first-centre-smooth-part}
\uses{mc:centres-avoid-open,mc:embedding-properties,mc:embedding-properties-smooth-centre,ana:nonkahler-locus,lag:regular-point-basics,ana:regular-transport,ana:strict-transform-sequence,ana:strict-transform-blowup,ana:regular-component,ana:analytic-preimage,rem:open-subset-analytic,hir:resolution-driver}
\source{wlodarczyk05-hironaka}
\lean{Hironaka.EmbeddedFromResolution.regularOpen,Hironaka.EmbeddedFromResolution.isClosed_singularLocus}\upstreamlean{TranscendentalNakamaye.Pending.EmbeddedFromResolution.mem_regularOpen_iff}\lean{Hironaka.EmbeddedFromResolution.regularAtDim_of_mem_regularOpen,Hironaka.EmbeddedFromResolution.mem_regularOpen_of_regularAtDim}\upstreamlean{TranscendentalNakamaye.Pending.EmbeddedFromResolution.strictTransform_eq_connectedComponentIn}\lean{Hironaka.EmbeddedFromResolution.mem_maxLocus_iff_of_centresAvoid}\upstreamlean{TranscendentalNakamaye.Pending.EmbeddedFromResolution.strictTransform_eq_connectedComponentIn_maxLocus}
\mixedleanok
Let \(X\) be compact, \(V\subset X\) a closed irreducible analytic subvariety and \(U=X\setminus V_{\mathrm{sing}}\), and suppose that the centres of \(T\) miss the
preimages of \(U\).
\begin{enumerate}
\item \(V_{\mathrm{sing}}\) is closed. If \(C\) is a closed submanifold of \(X_T\) which over \(U\) is the preimage of \(V\), and \(c\in C\) lies over \(U\), then
\(p_T(\operatorname{Exc}p_T)\subset V_{\mathrm{sing}}\) and the strict transform \(V_T\) of \(V\) is the connected component of \(C\) through \(c\); in particular
\(V_T\) is regular of dimension \(\dim V\) at each of its points.
\item Let \(\operatorname{inv}\) have the embedding properties in dimension \(\le k\), with \(\underline{\mathcal I}_V\) in its domain, and let \(T\) be a canonical
segment of \(\underline{\mathcal I}_V\). If the maximum locus \(C\) of \(\operatorname{inv}_{\underline{\mathcal I}_V}(T)\) on \(\operatorname{cosupp}(\underline{\mathcal I}_V)_T\)
contains a point \(c\) over \(U\), then over \(U\) it is the preimage of \(V\), and (1) applies to it.
\end{enumerate}
\end{lemma}
\begin{proof}
\mixedleanok
\(V_{\mathrm{sing}}\) is closed because regularity of a fixed dimension is an open condition on \(V\) (Lemma~\ref{lag:regular-point-basics}).
(1) By Lemma~\ref{mc:centres-avoid-open}, \(p_T(\operatorname{Exc}p_T)\subset X\setminus U=V_{\mathrm{sing}}\), and over the open set \(O=p_T^{-1}(U)\) the strict transform \(S=V_T\)
is \(p_T^{-1}(V)\). As \(p_T\) has bijective differential at \(c\) and \(V\) is regular of dimension \(\dim V\) at \(p_T(c)\in V_{\mathrm{reg}}\), \(S\) is regular of dimension
\(\dim V\) at \(c\) (Lemma~\ref{ana:regular-transport}), and so is \(C\), which agrees with \(S\) on \(O\) (Lemma~\ref{lag:regular-point-basics}). Since
\(V\not\subset p_T(\operatorname{Exc}p_T)\), \(S\) is a closed irreducible analytic subvariety (Lemma~\ref{ana:strict-transform-sequence}), hence the closure of the
connected component \(S_c\) of its regular locus through \(c\) (Lemma~\ref{ana:strict-transform-blowup}). \(C\) is analytic (Lemma~\ref{ana:analytic-preimage}) and all its
points are regular; let \(K\) be its connected component through \(c\). By the identity theorem along a component of a regular locus (an analytic set containing
a nonempty relatively open part of the component contains the component; proof of Proposition~\ref{ana:regular-component}), applied on \(O\ni c\) where \(S\) and
\(C\) agree, \(S_c\subset C\) and \(K\subset S\). The connected set \(S_c\) of regular points of \(C\) meets \(K\), so \(S_c\subset K\), and \(S=\overline{S_c}\subset C\).
Let \(x\in S\); then \(x\in C\) has a neighbourhood \(W\) with \(C\cap W\) inside the component \(K_x\) of \(C\) through \(x\)
(Proposition~\ref{ana:regular-component}(1)), and \(W\) meets \(S_c\subset K\), so \(K_x=K\) and \(x\in K\). Hence \(S=K\). All points of \(K\) are regular points of \(C\)
of dimension \(\dim V\) (Proposition~\ref{ana:regular-component}(1)), and near each of them \(C\) lies in \(K=S\subset C\), so \(S\) is regular of dimension \(\dim V\) there.
(2) \(T\) is a multiple test blow-up and \(C\) an admissible centre, in particular a closed submanifold (Corollary~\ref{mc:embedding-properties-smooth-centre}).
Over \(U\) the cosupport is the preimage of \(V\) (Lemma~\ref{mc:centres-avoid-open}(3)) and contains \(C\). \(V\cap U=V_{\mathrm{reg}}\) is a closed submanifold of \(U\) of pure
dimension \(\dim V\) (Lemma~\ref{rem:open-subset-analytic} for the traces), so by Corollary~\ref{mc:embedding-properties-smooth-centre} (the form of
Lemma~\ref{mc:smooth-part-invariant}) the invariant is constant on the cosupport over \(U\); as \(c\) lies in the maximum locus, so does every point of the cosupport
over \(U\).
\end{proof}

\begin{lemma}[The first centre meeting an open set]
\label{mc:first-centre-exists}
\uses{hir:canonical-segments,hir:resolution-driver,mc:centres-avoid-open,ana:blowup-surjective,hir:blowup-tower}
\source{wlodarczyk05-hironaka}
\lean{Hironaka.EmbeddedFromResolution.exists_first_centre,Hironaka.EmbeddedFromResolution.not_centresAvoid_of_cosupp_eq_empty}
\leanok
\begin{enumerate}
\item Let \(\operatorname{inv}\) be any value function, \(U\subset X\) open and \(T\) a canonical segment of a marked ideal on \(X\) (each centre empty or the
maximum locus of \(\operatorname{inv}\) of the preceding part on the cosupport) whose centres do not all miss the preimages of \(U\). Then \(T=T'S\) where \(T'\) is a
canonical segment whose centres miss the preimages of \(U\) and whose maximum locus, the next centre of \(T\), meets the preimage of \(U\).
\item Let \(X\) be compact, \(V\subset X\) a closed irreducible analytic subvariety with \(\underline{\mathcal I}_V\) in the domain of the invariants in
dimension \(\le k\), and \(T\) a multiple test blow-up of \(\underline{\mathcal I}_V\) with empty final cosupport. Then not all centres of \(T\) miss the preimages of
\(X\setminus V_{\mathrm{sing}}\).
\end{enumerate}
\end{lemma}
\begin{proof}
\leanok
(1) Induction on \(T\). Let \(T\) be \(T_0\) followed by the blowup along \(C\). If not all centres of \(T_0\) miss the preimages of \(U\), the induction hypothesis
applies to \(T_0\). Otherwise \(C\) meets the preimage of \(U\); so \(C\neq\emptyset\), \(C\) is the maximum locus at \(T_0\), and \(T'=T_0\).
(2) Otherwise, take \(x\in V_{\mathrm{reg}}\) (nonempty by definition of the subvariety); \(p_T\) is onto (Lemma~\ref{ana:blowup-surjective}), and a point over \(x\) lies in the
cosupport by Lemma~\ref{mc:centres-avoid-open}(3), against its emptiness.
\end{proof}

\subsection{Derivative ideals under a blowup}

In this subsection \(\pi\colon Y\to X\) is the blowup of the complex manifold \(X\) along \(C\)
(Definition~\ref{ana:blowups}), \(D\) its exceptional divisor, so that \(\mathcal O_Y(-qD)=\pi^*(\mathcal I_C^q)\)
(Lemma~\ref{bcalc:exceptional-divisor}), and \(\sigma^c(\mathfrak a,\mu)=(\pi^*\mathfrak a:\mathcal O_Y(-\mu D))\) is the
controlled transform (Definition~\ref{bcalc:controlled-transform-def}). Marked ideals are taken with \(N\) the whole
manifold. For a chart \(e\) and \(v\in\mathbf C^n\), \(\partial^e_v\) is the partial derivative of
Definition~\ref{hir:derivative-ideal}; \(e_z\) is a chart at \(z\) and \(e_1,\dots,e_n\) the standard basis.


\begin{definition}[Derivatives of a marked ideal]
\label{mc:derivative-marked}
\uses{hir:marked-ideal,hir:derivative-ideal,hir:cosupport}
\source{wlodarczyk05-hironaka}
\lean{Hironaka.DerivativeIdealTower.derivMarked,Hironaka.DerivativeIdealTower.derivMarked_N,Hironaka.DerivativeIdealTower.derivMarked_E,Hironaka.DerivativeIdealTower.derivMarked_ideal,Hironaka.DerivativeIdealTower.derivMarked_d,Hironaka.DerivativeIdealTower.cosupport_subset_of_le,Hironaka.DerivativeIdealTower.controlledTransform_le_of_le,Hironaka.DerivativeIdealTower.cosupp_subset_cosupp_derivMarked}
\leanok
For a marked ideal \(\underline{\mathcal I}=(M,N,E,\mathfrak a,d)\) and \(j\in\mathbb N\), \(\mathcal D^j(\underline{\mathcal I})=(M,N,E,\mathcal D^j(\mathfrak a),d-j)\).
If \(\mathfrak b\subset\mathfrak c\), then \(\operatorname{cosupp}(\mathfrak c,d)\subset\operatorname{cosupp}(\mathfrak b,d)\) and
\(\sigma^{\mathrm c}(\mathfrak b,\mu)\subset\sigma^{\mathrm c}(\mathfrak c,\mu)\). For \(N=M\) and \(j\le d\),
\(\operatorname{cosupp}\underline{\mathcal I}\subset\operatorname{cosupp}\mathcal D^j(\underline{\mathcal I})\).
\end{definition}
The two monotonicity statements are immediate from the definitions (a germ of \(\mathfrak b\) is a germ of \(\mathfrak c\); \(\sigma^*\mathfrak b\subset\sigma^*\mathfrak c\)), and the
last one is Lemma~\ref{hir:giraud-support} applied \(j\) times, \(\operatorname{cosupp}(\mathfrak a,(d-j)+j)\subset\operatorname{cosupp}(\mathcal D^j\mathfrak a,d-j)\).

\begin{lemma}[Derivative ideals along multiple test blow-ups]
\label{mc:inclusions-tower}
\uses{mc:derivative-marked,mc:inclusions,hir:geo-transform,hir:geo-transform-ambient,hir:bridge-consumers,hir:giraud-support,hir:derivative-ideal-generators}
\source{wlodarczyk05-hironaka}
\lean{Hironaka.DerivativeIdealTower.inclusions_tower}
\leanok
Let \(\underline{\mathcal I}=(M,M,E,\mathfrak a,d)\) with \(\mathfrak a\) coherent, \(j\le d\), and let \(T\) be a multiple test blow-up of \(\underline{\mathcal I}\)
(Definition~\ref{hir:geo-transform}). Then \(T\) is a multiple test blow-up of \(\mathcal D^j(\underline{\mathcal I})\), and
\[
 [\mathcal D^j(\underline{\mathcal I})]_T\subset\mathcal D^j(\underline{\mathcal I}_T),
\]
i.e.\ the transformed ideal of \(\mathcal D^j(\underline{\mathcal I})\) is contained in \(\mathcal D^j\) of the transformed ideal of \(\underline{\mathcal I}\); both transforms
have \(N\) the whole space, the same \(E\), and coherent ideals.
\end{lemma}
\begin{proof}
\leanok
Induction on \(T\); for the empty tower both sides are \(\mathcal D^j(\mathfrak a)\), coherent by Lemma~\ref{hir:derivative-ideal-generators}. Let \(T\) be followed by
the blowup \(\sigma\) along an admissible centre \(C\) of \(\underline{\mathcal I}_T=(M_T,M_T,E_T,\mathfrak a_T,d)\), and write
\(\mathfrak b_T\subset\mathcal D^j(\mathfrak a_T)\) for the transformed ideal of \(\mathcal D^j(\underline{\mathcal I})\), with mark \(d-j\) and the same \(E_T\).
\(C\) is admissible for \((M_T,M_T,E_T,\mathfrak b_T,d-j)\): it is a closed submanifold with normal crossings with \(E_T\), and
\(C\subset\operatorname{cosupp}(\mathfrak a_T,d)\subset\operatorname{cosupp}(\mathcal D^j\mathfrak a_T,d-j)\subset\operatorname{cosupp}(\mathfrak b_T,d-j)\)
(Definition~\ref{mc:derivative-marked}). By Lemma~\ref{hir:bridge-consumers}, \(\mathfrak a_T\subset\mathcal I_C^{d}\) and \(\mathfrak b_T\subset\mathcal I_C^{d-j}\), so both
transformed ideals are the coherent controlled transforms (Lemma~\ref{hir:geo-transform-ambient}), and by monotonicity and
Lemma~\ref{mc:inclusions} (with mark \(d-j\) and \(r=j\))
\[
 \sigma^{\mathrm c}(\mathfrak b_T,d-j)\subset\sigma^{\mathrm c}(\mathcal D^j\mathfrak a_T,d-j)\subset\mathcal D^j(\sigma^{\mathrm c}(\mathfrak a_T,d)).
\]
\(N\) stays the whole space and both \(E\) are transformed alike.
\end{proof}
This is Włodarczyk's Lemma ``le:inclusions2''. His Lemma ``le:Giraud'' (the cosupports of the transforms lie on the strict transforms of a
hypersurface of maximal contact) is Theorem~\ref{mc:giraud-tower} below: it follows from Lemma~\ref{mc:giraud-transform} and parts (2)--(3) of
``le:Gi'' on open sets (Lemma~\ref{bcalc:giraud-local}) by induction on the tower.

\subsection{Sums of marked ideals and the coefficient ideal}

In this subsection all marked ideals are \((M,M,E,\mathfrak a,d)\) with \(N=M\); for ideals \(\mathfrak J_0,\dots,\mathfrak J_{m-1}\),
\(\sum_{i<m}\mathfrak J_i\) is the ideal generated stalkwise by their germs.


\begin{lemma}[Finite sums of ideals]
\label{mc:finite-sum}
\uses{mc:homogenized-coherent,mc:coherent-germs,hir:controlled-operations,hir:cosupport-basic,bcalc:controlled-transform-def}
\lean{Hironaka.CoefficientIdeal.finSumIdeal_one,Hironaka.CoefficientIdeal.finSumIdeal_congr,Hironaka.CoefficientIdeal.finSumIdeal_le,Hironaka.CoefficientIdeal.isCoherent_finSumIdeal_of_lt,Hironaka.CoefficientIdeal.mem_cosupport_finSumIdeal,Hironaka.CoefficientIdeal.controlledTransform_finSumIdeal}
\leanok
Let \(\mathfrak J_0,\dots,\mathfrak J_{m-1}\) be coherent. Then \(\sum_{i<m}\mathfrak J_i\) is coherent; it is contained in every coherent ideal containing all
\(\mathfrak J_i\); \(\operatorname{cosupp}(\sum_{i<m}\mathfrak J_i,\mu)=\bigcap_{i<m}\operatorname{cosupp}(\mathfrak J_i,\mu)\); and if \(m\ge1\) and
\(\mathfrak J_i\subset\mathcal I_C^\mu\) for all \(i\), then \(\sigma^{\mathrm c}(\sum_{i<m}\mathfrak J_i,\mu)=\sum_{i<m}\sigma^{\mathrm c}(\mathfrak J_i,\mu)\).
\end{lemma}
\begin{proof}
\leanok
Coherence: Lemma~\ref{mc:homogenized-coherent}. A germ of the sum is a holomorphic combination of germs of the \(\mathfrak J_i\); this gives the
inclusion and, since a combination of germs vanishing to order \(\mu\) vanishes to order \(\mu\), the cosupport (the reverse inclusion because each
\(\mathfrak J_i\) is contained in the sum). The controlled transform: induction on \(m\), starting from \(\sum_{i<1}\mathfrak J_i=\mathfrak J_0\), with
\(\sum_{i<m+1}\mathfrak J_i=\sum_{i<m}\mathfrak J_i+\mathfrak J_m\) and \(\sigma^{\mathrm c}(\mathfrak a+\mathfrak b,\mu)=\sigma^{\mathrm c}(\mathfrak a,\mu)+\sigma^{\mathrm c}(\mathfrak b,\mu)\)
for coherent \(\mathfrak a,\mathfrak b\subset\mathcal I_C^\mu\) (Lemma~\ref{hir:controlled-operations}); the partial sums are contained in the coherent ideal \(\mathcal I_C^\mu\).
\end{proof}

\begin{definition}[Sum of finitely many marked ideals]
\label{mc:marked-finite-sum}
\uses{hir:marked-sum-product,mc:finite-sum,hir:marked-ideal-operations,hir:geo-transform-ambient,hir:bridge-consumers}
\source{wlodarczyk05-hironaka}
\lean{Hironaka.CoefficientIdeal.sumExp,Hironaka.CoefficientIdeal.mul_sumExp,Hironaka.CoefficientIdeal.one_le_sumExp,Hironaka.CoefficientIdeal.msum,Hironaka.CoefficientIdeal.IsAmbientFamily,Hironaka.CoefficientIdeal.IsAmbientFamily.isCoherent_msum,Hironaka.CoefficientIdeal.cosupp_msum,Hironaka.CoefficientIdeal.isAdmissibleCentre_msum_iff,Hironaka.CoefficientIdeal.geoTransform_msum,Hironaka.CoefficientIdeal.IsAmbientFamily.geoTransform}
\leanok
For marked ideals \(\underline{\mathcal I}_i=(M,N_i,E_i,\mathfrak a_i,\mu_i)\), \(i<m\),
\[
 \underline{\mathcal I}_0+\dots+\underline{\mathcal I}_{m-1}=\Bigl(M,N_0,E_0,\sum_{i<m}\mathfrak a_i^{\prod_{j\ne i}\mu_j},\prod_{j<m}\mu_j\Bigr).
\]
If \(m\ge1\), \(N_i=M\), \(E_i=E_0\), \(\mathfrak a_i\) is coherent and \(\mu_i\ge1\) for all \(i\), then the sum has coherent ideal and cosupport
\(\bigcap_i\operatorname{cosupp}\underline{\mathcal I}_i\); a centre is admissible for it if and only if it is admissible for every \(\underline{\mathcal I}_i\), and then the
transform of the sum is the sum of the transforms (Definition~\ref{hir:geo-transform}), whose summands again satisfy these hypotheses.
\end{definition}
\begin{proof}
\leanok
Write \(\mu=\prod_j\mu_j\) and \(e_i=\prod_{j\ne i}\mu_j\ge1\), so \(\mu_ie_i=\mu\). By Lemma~\ref{mc:finite-sum} and
\(\operatorname{cosupp}(\mathfrak a^{e},\nu e)=\operatorname{cosupp}(\mathfrak a,\nu)\) for \(e\ge1\) (Lemma~\ref{hir:marked-ideal-operations}),
\(\operatorname{cosupp}(\sum_i\mathfrak a_i^{e_i},\mu)=\bigcap_i\operatorname{cosupp}(\mathfrak a_i^{e_i},\mu_ie_i)=\bigcap_i\operatorname{cosupp}(\mathfrak a_i,\mu_i)\). The
admissibility statement follows, as all \(E_i\) agree. For an admissible centre \(C\), \(\mathfrak a_i\subset\mathcal I_C^{\mu_i}\) (Lemma~\ref{hir:bridge-consumers}), so
\(\mathfrak a_i^{e_i}\subset\mathcal I_C^{\mu}\), and by Lemmas~\ref{hir:geo-transform-ambient}, \ref{mc:finite-sum} and~\ref{hir:controlled-operations}
\(\sigma^{\mathrm c}(\sum_i\mathfrak a_i^{e_i},\mu)=\sum_i\sigma^{\mathrm c}(\mathfrak a_i^{e_i},\mu_ie_i)=\sum_i\sigma^{\mathrm c}(\mathfrak a_i,\mu_i)^{e_i}\); the strict transform of \(M\)
is the whole blowup, the \(E\) are transformed alike, the marks do not change, and the controlled transforms are coherent.
\end{proof}

\begin{lemma}[Multiple test blow-ups of a sum]
\label{mc:marked-finite-sum-tower}
\uses{mc:marked-finite-sum,hir:geo-transform}
\source{wlodarczyk05-hironaka}
\lean{Hironaka.CoefficientIdeal.geoTransformAlong_msum}
\leanok
Under the hypotheses of Definition~\ref{mc:marked-finite-sum}, the multiple test blow-ups of \(\underline{\mathcal I}_0+\dots+\underline{\mathcal I}_{m-1}\) are exactly
the simultaneous multiple test blow-ups of the \(\underline{\mathcal I}_i\), and along them the transform of the sum is the sum of the transforms.
\end{lemma}
\begin{proof}
\leanok
Induction on the tower, using at each step that the transforms satisfy the hypotheses again and Definition~\ref{mc:marked-finite-sum}.
\end{proof}
This is part (1) of Włodarczyk's Lemma ``le:operations'' for \(m\) summands; for \(m=2\) it is Lemma~\ref{hir:marked-operations-transform}(1).

\begin{definition}[Coefficient ideal]
\label{mc:coefficient-ideal}
\uses{mc:marked-finite-sum,mc:derivative-marked}
\source{wlodarczyk05-hironaka,wlodarczyk09-analytic-resolution:page-0022}
\lean{Hironaka.CoefficientIdeal.derivMarked_zero,Hironaka.CoefficientIdeal.coeffMarked,Hironaka.CoefficientIdeal.isAmbientFamily_derivMarked}
\leanok
The \emph{coefficient ideal} of \(\underline{\mathcal I}=(M,N,E,\mathfrak a,\mu)\) is
\[
 \mathcal C(\underline{\mathcal I})=\underline{\mathcal I}+\mathcal D(\underline{\mathcal I})+\dots+\mathcal D^{\mu-1}(\underline{\mathcal I}),
\]
i.e.\ \((\mathfrak a,\mu)+(\mathcal D\mathfrak a,\mu-1)+\dots+(\mathcal D^{\mu-1}\mathfrak a,1)\) (Definitions~\ref{mc:derivative-marked} and~\ref{mc:marked-finite-sum}).
For \(N=M\) and \(\mathfrak a\) coherent, the summands satisfy the hypotheses of Definition~\ref{mc:marked-finite-sum}.
\end{definition}

\begin{lemma}[The coefficient ideal is equivalent to the marked ideal]
\label{mc:coefficient-equivalent}
\uses{mc:coefficient-ideal,mc:marked-finite-sum-tower,mc:inclusions-tower,mc:derivative-marked}
\source{wlodarczyk05-hironaka,wlodarczyk09-analytic-resolution:page-0022}
\lean{Hironaka.CoefficientIdeal.coeff_tower}
\leanok
Let \(\underline{\mathcal I}=(M,M,E,\mathfrak a,\mu)\) with \(\mathfrak a\) coherent and \(\mu\ge1\). A tower is a multiple test blow-up of
\(\mathcal C(\underline{\mathcal I})\) if and only if it is one of \(\underline{\mathcal I}\), and along it
\(\operatorname{cosupp}\mathcal C(\underline{\mathcal I})_T=\operatorname{cosupp}\underline{\mathcal I}_T\).
\end{lemma}
\begin{proof}
\leanok
By Lemma~\ref{mc:marked-finite-sum-tower} the multiple test blow-ups of \(\mathcal C(\underline{\mathcal I})\) are the simultaneous ones of the
\(\mathcal D^i(\underline{\mathcal I})\), \(i<\mu\); the summand \(i=0\) is \(\underline{\mathcal I}\), and by Lemma~\ref{mc:inclusions-tower} every multiple test blow-up
of \(\underline{\mathcal I}\) is one of every \(\mathcal D^i(\underline{\mathcal I})\). Along such a tower,
\(\operatorname{cosupp}\mathcal C(\underline{\mathcal I})_T=\bigcap_{i<\mu}\operatorname{cosupp}[\mathcal D^i(\underline{\mathcal I})]_T\); the term \(i=0\) is
\(\operatorname{cosupp}\underline{\mathcal I}_T\), and for every \(i\), \(\operatorname{cosupp}\underline{\mathcal I}_T\subset\operatorname{cosupp}\mathcal D^i(\underline{\mathcal I}_T)\subset
\operatorname{cosupp}[\mathcal D^i(\underline{\mathcal I})]_T\) by Definition~\ref{mc:derivative-marked} and \([\mathcal D^i(\underline{\mathcal I})]_T\subset\mathcal D^i(\underline{\mathcal I}_T)\)
(Lemma~\ref{mc:inclusions-tower}).
\end{proof}
This is Włodarczyk's Lemma ``le:coeff'', \(\mathcal C(\underline{\mathcal I})\simeq\underline{\mathcal I}\); the source assumes maximal order, which the proof does not use.

\subsection{Restriction of the coefficient ideal to a submanifold}

In this subsection \(\underline{\mathcal B}=(M,M,E,\mathfrak b,d)\) has \(N=M\) and \(\mathfrak b\) coherent, and \(\underline{\mathcal I}=(M,M,E,\mathfrak a,\mu)\) has
\(N=M\), \(\mathfrak a\) coherent and \(\mu\ge1\). Sources: Włodarczyk's Lemmas ``le:cf'' and ``le:S''; in the analytic paper
(\texttt{wlodarczyk09-analytic-resolution}) Lemma~5.6.2, Lemma~5.6.3 and Proposition~5.6.4.


\begin{definition}[Restriction of a marked ideal to a set]
\label{mc:restrict-marked}
\uses{hir:marked-ideal,hir:order,mc:strict-along,mc:coefficient-ideal}
\source{wlodarczyk05-hironaka}
\lean{Hironaka.CoefficientIdealRestriction.withN,Hironaka.CoefficientIdealRestriction.withN_N,Hironaka.CoefficientIdealRestriction.withN_E,Hironaka.CoefficientIdealRestriction.withN_ideal,Hironaka.CoefficientIdealRestriction.withN_d,Hironaka.CoefficientIdealRestriction.cosupp_withN,Hironaka.CoefficientIdealRestriction.CentresIn,Hironaka.CoefficientIdealRestriction.centresIn_refl,Hironaka.CoefficientIdealRestriction.centresIn_step_iff,Hironaka.CoefficientIdealRestriction.N_geoTransformAlong_withN}\upstreamlean{TranscendentalNakamaye.Pending.CoefficientIdealRestriction.N_coeffMarked}\lean{Hironaka.CoefficientIdealRestriction.isCoherent_coeffMarked}
\mixedleanok
For a set \(S\subset M\), the \emph{restriction} \(\underline{\mathcal B}|_S=(M,S,E,\mathfrak b,d)\) has the same \(E\), ideal and mark and \(N=S\); its cosupport is
\(\{x\in S\mid \operatorname{ord}_x(\mathfrak b|_S)\ge d\}\) (Definition~\ref{hir:order}). A tower \(T\) has its \emph{centres in the strict transforms of \(S\)} if
each centre \(C_i\) lies in the strict transform \(S_i\) of \(S\) (Definition~\ref{mc:strict-along}). The coefficient ideal \(\mathcal C(\underline{\mathcal I})\)
has \(N=M\) and a coherent ideal.
\end{definition}

\begin{lemma}[Restriction along multiple test blow-ups]
\label{mc:restrict-tower}
\uses{mc:restrict-step,mc:restrict-order,mc:strict-along,hir:test-tower,hir:restricted-equivalence-basic,bcalc:strict-ideal}
\source{wlodarczyk05-hironaka}
\lean{Hironaka.CoefficientIdealRestriction.centresIn_of_isGeoTestTower_withN,Hironaka.CoefficientIdealRestriction.restrict_tower}
\leanok
Let \(S\subset M\) be a closed submanifold. A multiple test blow-up of \(\underline{\mathcal B}|_S\) has its centres in the strict transforms of \(S\).
Every multiple test blow-up \(T\) of \(\underline{\mathcal B}\) with centres in the strict transforms of \(S\) is a multiple test blow-up of
\(\underline{\mathcal B}|_S\), and along it \([\underline{\mathcal B}|_S]_T\) is restricted-equivalent to \((\underline{\mathcal B}_T)|_{S_T}\); moreover
\(\underline{\mathcal B}_T\) has \(N\) the whole space, both transformed ideals are coherent and \(S_T\) is a closed submanifold.
\end{lemma}
\begin{proof}
\leanok
An admissible centre lies in the cosupport, hence in \(N\), which is the strict transform of \(S\). For the second statement induct on \(T\). At a step
with centre \(C\subset S_T\) admissible for \(\underline{\mathcal B}_T\), \(C\) is admissible for \((\underline{\mathcal B}_T)|_{S_T}\) by
Lemma~\ref{mc:restrict-order}, hence for the restricted-equivalent \([\underline{\mathcal B}|_S]_T\); restricted-equivalent marked ideals have the same
transform (Lemma~\ref{hir:restricted-equivalence-basic}), and Lemma~\ref{mc:restrict-step} applies. The strict transform of a closed submanifold
containing the centre is a closed submanifold (Lemma~\ref{bcalc:strict-ideal}).
\end{proof}
This is the last statement of Włodarczyk's Lemma ``le:cf'': \([(\mathcal I,\mu)|_S]_i=(\mathcal I_i,\mu)|_{S_i}\).

\begin{lemma}[The support of the restricted coefficient ideal]
\label{mc:coefficient-restriction-support}
\uses{mc:restrict-marked,mc:restrict-order,mc:coefficient-equivalent,mc:finite-sum,hir:order-power,hir:giraud-support}
\source{wlodarczyk05-hironaka,wlodarczyk09-analytic-resolution:page-0023}
\lean{Hironaka.CoefficientIdealRestriction.cosupport_coeff_subset,Hironaka.CoefficientIdealRestriction.cosupp_coeff_withN}
\leanok
For every set \(S\subset M\), \(\operatorname{cosupp}\underline{\mathcal I}\cap S=\operatorname{cosupp}(\mathcal C(\underline{\mathcal I})|_S)\).
\end{lemma}
\begin{proof}
\leanok
The inclusion \(\subset\) is Lemma~\ref{mc:coefficient-equivalent} (for the empty tower) followed by Lemma~\ref{mc:restrict-order}. Conversely, let
\(x\in S\) with \(\operatorname{ord}_x(\mathcal C(\mathfrak a)|_S)\ge\mu_{\mathcal C}\). The ideal of \(\mathcal C(\underline{\mathcal I})\) contains
\(g^e\) for every \(g\in\mathcal D^{\mu-1}\mathfrak a\), where \(e\) is the exponent of the last summand \((\mathcal D^{\mu-1}\mathfrak a,1)\), and
\(\mu_{\mathcal C}=1\cdot e\ge1\). So \(g^e\), hence \(g\), vanishes at \(x\): \(x\in V(\mathcal D^{\mu-1}\mathfrak a)=\operatorname{cosupp}(\mathfrak a,\mu)\)
(Lemma~\ref{hir:giraud-support}).
\end{proof}
This is the first statement of Włodarczyk's Lemma ``le:S'' (Proposition~5.6.4 of the analytic paper), with \(\mu_{\mathcal C}=\mu!\) in place of
\(\operatorname{lcm}(1,\dots,\mu)\), which changes nothing. The source proves \(\supset\) by expanding the germs of \(\mathfrak a\) in coordinates normal to \(S\);
the last summand of \(\mathcal C\) makes this unnecessary, and \(S\) need not be a submanifold.

\begin{lemma}[The restricted coefficient ideal along multiple test blow-ups]
\label{mc:coefficient-restriction-tower}
\uses{mc:restrict-tower,mc:coefficient-equivalent,mc:restrict-order,mc:restrict-marked}
\source{wlodarczyk05-hironaka,wlodarczyk09-analytic-resolution:page-0023}
\lean{Hironaka.CoefficientIdealRestriction.coeff_restrict_tower}
\leanok
Let \(S\subset M\) be a closed submanifold and \(T\) a multiple test blow-up of \(\underline{\mathcal I}\) with centres in the strict transforms of \(S\).
Then \(T\) is a multiple test blow-up of \(\mathcal C(\underline{\mathcal I})|_S\), along it \([\mathcal C(\underline{\mathcal I})|_S]_T\) is
restricted-equivalent to \(\mathcal C(\underline{\mathcal I})_T|_{S_T}\), and
\(\operatorname{cosupp}\underline{\mathcal I}_T\cap S_T\subset\operatorname{cosupp}[\mathcal C(\underline{\mathcal I})|_S]_T\).
\end{lemma}
\begin{proof}
\leanok
By Lemma~\ref{mc:coefficient-equivalent}, \(T\) is a multiple test blow-up of \(\mathcal C(\underline{\mathcal I})\) and
\(\operatorname{cosupp}\mathcal C(\underline{\mathcal I})_T=\operatorname{cosupp}\underline{\mathcal I}_T\). Lemma~\ref{mc:restrict-tower} for
\(\underline{\mathcal B}=\mathcal C(\underline{\mathcal I})\) gives the first two statements, and the inclusion follows from Lemma~\ref{mc:restrict-order}.
\end{proof}
This is part (1) and the inclusion \(\subset\) of part (2) of Włodarczyk's Lemma ``le:S'' (Proposition~5.6.4 of the analytic paper). The reverse
inclusion \(\operatorname{cosupp}[\mathcal C(\underline{\mathcal I})|_S]_T\subset\operatorname{cosupp}\underline{\mathcal I}_T\) for a nonempty tower, and
with it part (3), is Lemma~\ref{mc:coefficient-restriction-full} below (the source's tracking of the Taylor coefficients in coordinates normal to the
\(S_i\) through the monomial charts).

\subsection{The restricted coefficient ideal along multiple test blow-ups: the full statement}

This subsection proves the remaining inclusion of Włodarczyk's Lemma ``le:S'' (Proposition~5.6.4 of the analytic paper), following the source:
the Taylor coefficients of the transforms in the normal coordinates of the strict transforms \(S_i\) of \(S\) are controlled transforms of
derivatives. Normal derivatives play the role of the coefficients \(c_\alpha=(\alpha!)^{-1}\partial^\alpha_xf|_S\).


\begin{definition}[Flat charts and normal derivatives]
\label{mc:flat-chart}
\uses{mc:normal-taylor-order,fpgeom:complex-analytic-subvariety}
\source{wlodarczyk05-hironaka}
\lean{Hironaka.NormalDerivativeChart.FlatChart,Hironaka.NormalDerivativeChart.mem_congr_of_isCoherent,Hironaka.NormalDerivativeChart.FlatChart.ψ,Hironaka.NormalDerivativeChart.FlatChart.inv,Hironaka.NormalDerivativeChart.FlatChart.nd,Hironaka.NormalDerivativeChart.FlatChart.ofRegularAtDim,Hironaka.NormalDerivativeChart.FlatChart.mem_flatDomain,Hironaka.NormalDerivativeChart.FlatChart.ψ_self,Hironaka.NormalDerivativeChart.FlatChart.zero_mem_target,Hironaka.NormalDerivativeChart.FlatChart.inv_zero,Hironaka.NormalDerivativeChart.FlatChart.eventually_mem_flatDomain,Hironaka.NormalDerivativeChart.FlatChart.mdifferentiableAt_ψ,Hironaka.NormalDerivativeChart.FlatChart.eventually_mdifferentiableAt_ψ,Hironaka.NormalDerivativeChart.FlatChart.mdifferentiableAt_inv,Hironaka.NormalDerivativeChart.FlatChart.eventually_mdifferentiableAt_inv,Hironaka.NormalDerivativeChart.FlatChart.eventually_inv_ψ,Hironaka.NormalDerivativeChart.FlatChart.eventually_ψ_inv,Hironaka.NormalDerivativeChart.FlatChart.eventually_flat,Hironaka.NormalDerivativeChart.FlatChart.tendsto_inv,Hironaka.NormalDerivativeChart.FlatChart.tendsto_ψ,Hironaka.NormalDerivativeChart.FlatChart.eventually_differentiableAt_comp_inv,Hironaka.NormalDerivativeChart.FlatChart.nd_nil,Hironaka.NormalDerivativeChart.FlatChart.mem_idealPow_of_vanishesToOrderOn,Hironaka.NormalDerivativeChart.FlatChart.injective_fderiv_Φ}
\leanok
A \emph{flat chart} for \(S\) at \(p\) is the preferred chart followed by a local biholomorphism \(\Phi\) with \(\Phi(e\,p)=0\) in which \(S\) is the plane
\(\mathbb C^k\times0\) near \(p\) (the data of a regular point). The \emph{normal derivatives} of \(f\) are \(\partial^L(f\circ\psi^{-1})\circ\psi\),
\(\psi=\Phi\circ e\). Lemma~\ref{mc:normal-taylor-order} applies to flat charts.
\end{definition}

\begin{lemma}[Normal derivatives of an ideal]
\label{mc:normal-derivative-ideal}
\uses{mc:flat-chart,hir:derivative-ideal,mc:coordinate-derivatives}
\source{wlodarczyk05-hironaka}
\lean{Hironaka.NormalDerivativeChart.FlatChart.nd_mem_iterDerivIdeal}
\leanok
For a coherent \(\mathfrak b\) and \(g\in\mathfrak b_p\), every derivative of \(g\) in a flat chart along a list \(L\) lies in \(\mathcal D^{|L|}(\mathfrak b)_p\).
\end{lemma}
\begin{proof}
\leanok
A coordinate derivative is a holomorphic combination of chart derivatives (Lemma~\ref{mc:coordinate-derivatives}); induct on \(L\).
\end{proof}

\begin{lemma}[The coefficient invariant]
\label{mc:normal-coefficient-invariant}
\uses{mc:flat-chart,mc:normal-derivative-ideal,mc:normal-taylor-order,mc:restrict-marked,mc:derivative-marked,ana:local-holomorphic-inverse}
\source{wlodarczyk05-hironaka}
\lean{Hironaka.CoefficientRestrictionTower.vanishesToOrderOn_of_isInSpanOnAt,Hironaka.CoefficientRestrictionTower.GoodAt,Hironaka.CoefficientRestrictionTower.GoodAt.orderGE,Hironaka.CoefficientRestrictionTower.goodAt_refl,Hironaka.CoefficientRestrictionTower.isInSpanAt_comp_map,Hironaka.CoefficientRestrictionTower.FlatChart.exists_of_map,Hironaka.CoefficientRestrictionTower.FlatChart.eventually_map_inv,Hironaka.CoefficientRestrictionTower.FlatChart.holomorphicAt_nd}
\leanok
Let \(\underline{\mathcal I}=(X,X,E,\mathfrak a,\mu)\), \(\mathfrak a\) coherent, \(S\) a closed submanifold and \(T\) a tower. A point \(p\in S_T\) is \emph{good} if there are a flat chart
for \(S_T\) at \(p\) and generators of \((\mathfrak a_T)_p\) whose normal derivatives \(\partial^L\), \(|L|<\mu\), restrict on \(S_T\) near \(p\) into
\([\mathcal D^{|L|}\mathfrak a]_T\). At a good point where every \([\mathcal D^i\mathfrak a]_T\) has order \(\ge\mu-i\) along \(S_T\), every germ of \(\mathfrak a_T\) lies in
\(\mathfrak m_p^\mu\). Every point of \(S\) is good for the empty tower.
\end{lemma}
\begin{proof}
\leanok
The first statement is Lemma~\ref{mc:normal-taylor-order} for the generators. For the empty tower take all germs of \(\mathfrak a\) as generators and
Lemma~\ref{mc:normal-derivative-ideal}.
\end{proof}

\begin{lemma}[The coefficient invariant along multiple test blow-ups]
\label{mc:normal-coefficient-step}
\uses{mc:normal-coefficient-invariant,mc:nested-product-chart,mc:normal-derivative-monomial,mc:monomial-local-equation,mc:restrict-tower,mc:inclusions-tower,bcalc:strict-submanifold,bcalc:monomial-lift,bcalc:isomorphism-off-centre,bcalc:function-transform,hir:geo-transform-ambient,hir:bridge-consumers,hir:geo-transform-wellformed,ana:restricted-colon}
\source{wlodarczyk05-hironaka}
\lean{Hironaka.CoefficientRestrictionTower.goodAt_step_of_notMem,Hironaka.CoefficientRestrictionTower.nd_monomial_identity,Hironaka.CoefficientRestrictionTower.goodAt_step_of_mem,Hironaka.CoefficientRestrictionTower.goodAt_tower}
\leanok
Along a multiple test blow-up of \(\underline{\mathcal I}\) with centres in the strict transforms of \(S\), every point of \(S_T\) is good.
\end{lemma}
\begin{proof}
\leanok
Induct on the tower; let \(\pi\) blow up \(C\subset S_T\) and \(z\in S_{T'}\). If \(\pi z\notin C\), \(\pi\) is a local biholomorphism near \(z\): pull back the chart and the
generators. If \(\pi z\in C\), refine the chart at \(\pi z\) to a nested product-form chart (Lemma~\ref{mc:nested-product-chart}) and take the monomial chart \(W\) at
\(z\) (Lemma~\ref{bcalc:strict-submanifold}): \(\sigma_j\circ W\) are the old coordinates, \(S_{T'}=W^{-1}(\mathbb C^k\times0)\), and \(W_j\) is a local equation of the
exceptional divisor (Lemma~\ref{mc:monomial-local-equation}). The new chart is \(W-W(z)\); the new generators are \(h'\) with \(g\circ\pi=W_j^\mu h'\) (they generate the
controlled transform after cancelling \(W_j^\mu\)). Lemma~\ref{mc:normal-derivative-monomial} gives
\(W_j^\mu\,\partial^Lh'=W_j^{|L|}(\partial^Lg)\circ\pi\); on \(S_T\) the right side is \(W_j^{|L|}\sum a_i(s_i\circ\pi)\) with \(s_i\in[\mathcal D^{|L|}\mathfrak a]_T\subset\mathcal I_C^{\mu-|L|}\),
so \(s_i\circ\pi=W_j^{\mu-|L|}s_i'\) with \(s_i'\in[\mathcal D^{|L|}\mathfrak a]_{T'}\); cancel \(W_j^\mu\) along \(S_{T'}\), where it vanishes on no relatively open set
(\(S_{T'}\) is the closure of \(\pi^{-1}(S_T\setminus C)\)).
\end{proof}

\begin{lemma}[The restricted coefficient ideal along multiple test blow-ups: full statement]
\label{mc:coefficient-restriction-full}
\uses{mc:normal-coefficient-step,mc:coefficient-restriction-tower,mc:coefficient-restriction-support,mc:marked-finite-sum-tower,mc:finite-sum,hir:order-power-along,hir:restricted-equivalence-basic}
\source{wlodarczyk05-hironaka,wlodarczyk09-analytic-resolution:page-0023}
\lean{Hironaka.CoefficientRestrictionTower.cosupp_coeff_withN_tower_subset,Hironaka.CoefficientRestrictionTower.cosupp_coeff_withN_tower,Hironaka.CoefficientRestrictionTower.coeff_restrict_tower_iff}
\leanok
Let \(\underline{\mathcal I}=(X,X,E,\mathfrak a,\mu)\) with \(\mathfrak a\) coherent and \(\mu\ge1\), and \(S\) a closed submanifold. A tower is a multiple test blow-up of
\(\mathcal C(\underline{\mathcal I})|_S\) if and only if it is a multiple test blow-up of \(\underline{\mathcal I}\) with centres in the strict transforms of \(S\), and along it
\(\operatorname{cosupp}[\mathcal C(\underline{\mathcal I})|_S]_T=\operatorname{cosupp}\underline{\mathcal I}_T\cap S_T\).
\end{lemma}
\begin{proof}
\leanok
The inclusion \(\supset\) is Lemma~\ref{mc:coefficient-restriction-tower}. Conversely, at \(x\) in the left side, \([\mathcal C(\underline{\mathcal I})|_S]_T\) is restricted-equivalent
to \(\mathcal C(\underline{\mathcal I})_T|_{S_T}=\sum_i[\mathcal D^i\mathfrak a]_T^{e_i}\) (Lemma~\ref{mc:marked-finite-sum-tower}); along \(S_T\) the order of the sum is the minimum of the
scaled orders (Lemma~\ref{hir:order-power-along}), so \([\mathcal D^i\mathfrak a]_T\) has order \(\ge\mu-i\) along \(S_T\) at \(x\), and the good point \(x\)
(Lemma~\ref{mc:normal-coefficient-step}) lies in \(\operatorname{cosupp}\underline{\mathcal I}_T\). The test blow-ups correspond by induction on the tower: a centre inside \(S_T\) is admissible
for one side iff for the other, since the cosupports correspond and \(E\) is the same.
\end{proof}
This is Włodarczyk's Lemma ``le:S'' (2)--(3) for the restriction keeping \(E\) (no member of \(E\) is dropped).
